\chapter{Үндсэн тохиргооны файлууд}
\label{AppendixB}
\pagecolor{white}

Энэ хавсралтад системийн байршуулалт болон автоматжуулалттай холбоотой үндсэн тохиргооны файлуудыг оруулав.

%-------------------------------------------------------------------------------

\section*{Б.1 Dockerfile — Node.js/Express бэкэнд}

\begin{lstlisting}[language={}, caption={backend/Dockerfile}]
FROM node:18-alpine
WORKDIR /app
COPY package*.json ./
RUN npm ci --only=production
COPY . .
EXPOSE 3000
CMD ["node", "server.js"]
\end{lstlisting}

Node.js 18 LTS-ийн alpine суурь дүрс ашигласнаар Docker дүрсний хэмжээ хамгийн бага байна. \code{npm ci --only=production} нь зөвхөн production хамаарлуудыг суулгаж, хөгжүүлэлтийн хамаарлуудыг оруулахгүй байлгана.

%-------------------------------------------------------------------------------

\section*{Б.2 GitHub Actions CI/CD Workflow}

\begin{lstlisting}[language={}, caption={.github/workflows/deploy.yml}]
name: Deploy to AWS EC2

on:
  push:
    branches: [main]
    paths:
      - 'backend/**'

jobs:
  deploy:
    runs-on: ubuntu-latest
    steps:
      - uses: actions/checkout@v3

      - name: Configure AWS credentials
        uses: aws-actions/configure-aws-credentials@v2
        with:
          aws-access-key-id: ${{ secrets.AWS_ACCESS_KEY_ID }}
          aws-secret-access-key: ${{ secrets.AWS_SECRET_ACCESS_KEY }}
          aws-region: ap-northeast-1

      - name: Login to Amazon ECR
        id: login-ecr
        uses: aws-actions/amazon-ecr-login@v1

      - name: Build and push Docker image
        env:
          ECR_REGISTRY: ${{ steps.login-ecr.outputs.registry }}
          IMAGE_TAG: ${{ github.sha }}
        run: |
          docker build \
            -t $ECR_REGISTRY/dz-zaal-backend:$IMAGE_TAG \
            -t $ECR_REGISTRY/dz-zaal-backend:latest \
            ./backend
          docker push $ECR_REGISTRY/dz-zaal-backend:$IMAGE_TAG
          docker push $ECR_REGISTRY/dz-zaal-backend:latest

      - name: Deploy to EC2
        uses: appleboy/ssh-action@master
        with:
          host: ${{ secrets.EC2_HOST }}
          username: ubuntu
          key: ${{ secrets.EC2_SSH_KEY }}
          script: |
            aws ecr get-login-password --region ap-northeast-1 | \
              docker login --username AWS --password-stdin \
              ${{ secrets.ECR_REGISTRY }}
            docker pull ${{ secrets.ECR_REGISTRY }}/dz-zaal-backend:latest
            docker stop dz-zaal-backend || true
            docker rm dz-zaal-backend || true
            docker run -d --name dz-zaal-backend \
              -p 3000:3000 \
              --env-file /home/ubuntu/.env \
              ${{ secrets.ECR_REGISTRY }}/dz-zaal-backend:latest
\end{lstlisting}

Workflow нь \code{backend/} хавтаст өөрчлөлт push хийх үед автоматаар ажиллана. AWS нэвтрэлтийн мэдээлэл (access key, SSH түлхүүр) нь GitHub Secrets-д хадгалагдаж кодод задрахгүй.

%-------------------------------------------------------------------------------

\section*{Б.3 Firebase Firestore Security Rules}

\begin{lstlisting}[language={}, caption={firestore.rules}]
rules_version = '2';
service cloud.firestore {
  match /databases/{database}/documents {

    match /users/{userId} {
      allow read, write: if request.auth != null
                         && request.auth.uid == userId;
    }

    match /bookings/{bookingId} {
      allow read: if request.auth != null
                  && resource.data.userId == request.auth.uid;
      allow create: if request.auth != null;
      allow update, delete: if request.auth != null
                             && resource.data.userId == request.auth.uid;
    }

    match /fixed_bookings/{docId} {
      allow read: if request.auth != null;
      allow write: if false;
    }

    match /email_verifications/{docId} {
      allow read, write: if false;
    }
  }
}
\end{lstlisting}

Firestore аюулгүй байдлын дүрэм нь хэрэглэгч зөвхөн өөрийн захиалгыг унших, засах эрхтэй болохыг баталгаажуулна. \code{email\_verifications} болон \code{fixed\_bookings} цуглуулгуудад клиент тал бичих боломжгүй.
